%% This document created by Scientific Word (R) Version 2.0
%% Starting shell: book
	

\documentclass[12pt,thmsa]{book}
\usepackage{amsfonts}
%%%%%
\usepackage{sw20lbck}

%TCIDATA{TCIstyle=Book/book.lat,lbck,book}

\input tcilatex
\QQQ{Language}{
American English
}

\begin{document}

\author{David Joyner}
\title{An application of Arthur's invariant trace formula for $SL(2)$}
\date{4-6-97}
\maketitle

\begin{abstract}
We apply the invariant trace formula to the study of cuspidal automorphic
representations of $SL(2,\Bbb{A})$ as compared with the genuine cuspidal
automorphic representations of the 2-fold metaplectic cover $\overline{SL(2,%
\Bbb{A})}$. The main tools are the trace formula, combined with results of
Arthur, Gelbart-Sally, Labesse-Langlands, and Waldspurger.
\end{abstract}

\tableofcontents

\chapter{Introduction}

The goals of this book are two-fold: on one hand this is a detailed
exposition of Arthur's invariant trace formula in the simple case of $SL(2)$
which may serve as an introduction to his work, and on the other hand, the
problem we have applied it to (in chapter 10) will hopefully be of some
interest to researchers in the field of the theory of automorphic
representations. Interestingly enough, our intuition regarding automorphic
representations of reductive groups gained from the Langlands philosophy
fails to offer proper guidance for automorphic representations of
metaplectic covers of reductive groups. This was one motivation for the
amount of detail included. The motivation for researching the problem
discussed in chapter 10 has its origins in several stimulating conversations
with R. Langlands in the late 1980's.

To possibly save the reader some time, we will briefly explain the purpose
of each chapter. Most chapters have an introduction so this will merely be
an overview.

Since we have in mind not only the study of the automorphic representations
of $SL(2)$ but also its 2-fold cover, we have put all the background
information on the metaplectic covering groups into a single chapter. As a
guiding principle in all of the chapters, when it is no harder to work with
n-fold covers we shall do so.

Since our techniques ultimately rely on the trace formula, a generalization
of Poisson summation in abelian harmonic analysis, it is not surprising that
several instances of Poisson summation must be used at various stages. We
present in the next chapter a detailed summary of all the cases of Poisson
summation which we will use. In addition, to give a feel for the argument in
chapter 10, we give in this chapter the analogous metaplectic correspondence
in the case where $SL(2)$ has been replaced by $GL(1)$. In addition, we
define some distributions which will be needed in an inductive step in the
derivation of the invariant trace formula in chapters 8 and 9. The fact that
the metaplectic $GL(1)$ comparison used these distributions gives us added
confidence that we have gotten them right.

Arthur's original derivation of the invariant trace formula required making
certain assumptions in local harmonic analysis. Since his original paper 
\cite{A6} appeared, these assumptions have been verified. The purpose of the
next chapter is to survey the representation theory of $SL(2)$ over a p-adic
field and verify directly some of these assumptions. For example, the
structure of the unitary dual of $SL(2)$ is explicitly recalled.
From this explicit information we can deduce the metaplectic analog
of the trace Paley-Wiener theorem \cite{BDK} for
the $2$-fold cover of $SL(2,{\Bbb{Q}}_p)$.

Chapter 5 presents the local correspondences needed in the comparison
between the invariant trace formula of $SL(2)$ and its 2-fold cover. The
results are rather technical to state but relatively easy to prove (at least
once you know what you want) compared to similar works in the field \cite{L4}%
, \cite{Fl}, \cite{R}. On the other hand, these comparisons are more
involved than those of \cite{JL}, \cite{LL}.

The chapter after that is the metaplectic analog of chapter 4 for the n-fold
cover of $SL(2)$. Unfortunately, much less is known in this case and we must
resort to stating conjectures. Fortunately, what we need for the derivation
of the invariant metaplectic trace formula in chapter 9 in the case of $n=2$
is known.

In the next chapter we use a construction of supercuspidals due to Kutzko
and Manderscheid along with the Frobenius formula to prove a local
correspondence between the supercuspidals of $SL(2)$ and its 2-fold cover.
This chapter owes much to discussions with Jason Adams and Jason
Schultz who pointed out an error in an earlier version.
This correspondence will be useful in chapter 10.

Chapter 8 derives the invariant trace formula for $SL(2)$. In fact, we go
further than is needed for the comparison in chapter 10: we derive the
stable trace formula from the invariant trace formula. Note that this is not
exactly the same as the stable trace formula of \cite{LL} since the
invariant trace formula gives a formula for the trace of the right regular
representation $R_{cusp}$ of $SL(2,\Bbb{A})$ on the space of
square-integrable cusp forms whereas the trace formula of \cite{LL} gives a
formula for the trace of the right regular representation of $SL(2,\Bbb{A})$
on the space of all square-integrable automorphic forms which occur
discretely in the spectrum. Note that (for simplicity) this derivation is
only for $SL(2)/{\Bbb{Q}}$ even though we shall actually use it for $SL(2)$
over any number field. This chapter may also serve as an introduction to the
papers \cite{LL}, \cite{A2}, \cite{A3}, \cite{A6}, \cite{A9}, \cite{A10}.

At this point we must confess that it would have been possible to establish
the comparisons of chapter 10 using the original (non-invariant) trace
formula of $SL(2)$ and its 2-fold cover. This wasn't discovered this until
after the material in chapters 8, 9 had already been written. This material
on the invariant trace formula and on the stable trace formula for $SL(2)$
serves as an exposition of Arthur's work and the work of Labesse-Langlands
and complements other introductions (\cite{Sho}, \cite{A5}, \cite{A15}, \cite
{Lab} ...). In any case, readers may skip any parts of chapters 8, 9 dealing
with the invariant trace formula or the stable trace formula, and may skip
the parts in chapters 4, 6 which establish the local assumptions for the
invariant trace formula (which means everything except for the character
formulas).

Chapter 9 is the metaplectic analog of most of chapter 8. We do not put the
spectral side of the metaplectic trace formula into stable form, only the
geometric side. This is almost forced on us because the structure of the
unitary dual of the n-fold cover of $SL(2)$ depends so much on the parity of 
$n$ and a case-by-case analysis would take us too far afield. Again, for
simplicity this derivation is only for the two-fold cover of
$SL(2,{\Bbb{A}}_{\Bbb{Q}})$ even though we shall
actually use it for the two-fold cover of $SL(2)$ over any number field.
The metaplectic analog of one of
Arthur's assumptions \cite{A6} in his derivation of the
invariant trace formula, the metaplectic analog of the
trace Paley-Wiener theorems in the $2$-adic case (\cite{CD}, \cite{Bou}, \cite{Ao})
is, as far as I know, unknown at this time.
In general, we must therefore assume that the
metaplectic analog of the Paley-Wiener theorem is true for the
$2$-fold cover of $SL(2,F)$, where $F$ is $2$-adic,
in our derivation of the invariant trace formula for
the 2-fold cover of $SL(2)$. 
In any case, as we mentioned in the previous paragraph, the main result
will not depend on the validity of this assumption. 

Chapter 10 is the application. Using the results of chapters 4, 5, 6, 7, 8
and 9, we compare the corresponding terms in the invariant
trace formula of $SL(2)$ and its 2-fold cover. The comparisons
yield an equality of the
following sort:
$\mathrm{tr}\;R_{cusp}(f)=\mathrm{tr}\;\overline{R}_{cusp}(\overline{f})$,
for all corresponding functions $f,\overline{f}$ in their
respective Hecke algebras. This, combined with some results of Waldspurger,
implies a relationship between the automorphic cuspidal
representations of $SL(2,\Bbb{A})$ and the genuine automorphic
cuspidal representations of the
2-fold cover $\overline{SL(2,\Bbb{A})}$. From this we obtain that each
``good'' cuspidal representation of $SL(2,{\Bbb{A}})$ occurs with
multiplicity one. Here ``good'' simply means that if $\pi =\otimes \pi _v$
and $v|2$ then $\pi _v$ is either a unitary principal series (in the
``continuous series'') or a complementary series. This is forced on us
because of our poor understanding of character relations of the discrete
series representations when $p=2$. This type of result has been anticipated
by \cite{Lab} and in some unpublished work of Y. Flicker. D. Blasius has
shown that the analog of this multiplicity one result is false for $SL(n),\
n>2$ \cite{Blas}.

We say nothing here about L-packets. However, we do note that the ``image of
this correspondence from $\overline{SL(2,\Bbb{A})}$ to $SL(2,\Bbb{A})$''
does not consist of the cuspidal automorphic representations with trivial
central character. This might be compared with Flicker's correspondence \cite
{Fl}, where the image has an ``even'' central character.

Finally, this was written on a sabbatical visiting the University of
Maryland, College Park. I'm grateful to the Naval Academy for giving me the
time off to think about these questions and
to the Mathematics Department at the
University of Maryland for the stimulating atmosphere.

\input book1.tex \input book2.tex \input book2a.tex \input book3.tex 
\input book3b.tex \input book4.tex \input book5nw.tex \input book6a.tex 
\input book6b.tex \input book6c.tex \input book6d.tex \input book6e.tex 
\input book6f.tex \input book6g.tex \input book7a.tex \input book7c.tex 
\input book7d.tex \input book7e.tex \input book7f.tex \input book7gnw.tex 
\input book7h.tex \input book8.tex \input bookrefn.tex

\end{document}
